\subsubsection{Fibrations over a projective base}

\begin{lemma}[Proved: a uniform bound for residual inseparability]
\label{D:residual-inseparability-bound}
Let $M/F$ be a finitely generated extension of fields of
characteristic $p>0$.  There is a finite purely inseparable
extension $F_0/F$, depending only on $M/F$, with the following
property: for every finite extension $L/F$ in a common
overfield for which $LM/M$ is separable, we have
\[
                 [L:F]_{\mathrm{insep}}\mid [F_0:F].
\]
We may choose any $F_0$ for which $MF_0/F_0$ is separable.
In characteristic zero the corresponding bound is one.
\end{lemma}
\begin{proof}
Choose an integral algebraic $F$-scheme with function field
$M$.  Lemma 33.27.4, tag 0BXT, supplies a finite purely
inseparable $F_0/F$ for which its normalization after base
change is geometrically normal
\cite{C_StacksLogFibration}.  Its function field $MF_0$
is therefore separable over $F_0$.

Put $M_0=MF_0$ and $L_0=LF_0$.  Since $LM/M$ is separable,
its base change $L_0M_0/M_0$ is separable.  A separable
extension $M_0/F_0$ cannot make an inseparable finite
algebraic extension separable.  To verify this assertion,
let $L_0^{\mathrm{sep}}$ be the maximal separable
subextension of $L_0/F_0$.  Then
$L_0M_0/L_0^{\mathrm{sep}}M_0$ is both purely inseparable
and separable, so is trivial.  Thus
$L_0\subset L_0^{\mathrm{sep}}M_0$.  The latter is separable
over $L_0^{\mathrm{sep}}$, so its purely inseparable finite
subextension $L_0/L_0^{\mathrm{sep}}$ is trivial.
Consequently $LF_0/F_0$ is separable.

Now let $L^{\mathrm{sep}}$ be the maximal separable
subextension of $L/F$.  The extension
$LF_0/L^{\mathrm{sep}}F_0$ is purely inseparable and,
by the preceding paragraph, separable, so it too is
trivial.  Hence $L\subset L^{\mathrm{sep}}F_0$ and
\[
 [L:F]_{\mathrm{insep}}
  =[L:L^{\mathrm{sep}}]
  \leq [L^{\mathrm{sep}}F_0:L^{\mathrm{sep}}]
  \leq [F_0:F].
\]
Both sides of the asserted divisibility are powers of
$p$, so this inequality proves it.
\end{proof}

\begin{proposition}[Proved: generic-fibre generation over a projective base]
\label{D:projective-base-generation}
Let $k$ be algebraically closed, let $f:X\to B$ be a
dominant morphism of normal connected proper $k$-schemes,
and suppose $B$ is projective and the geometric generic
fibre $X_{\overline\eta}$ is integral.  Let
$Y\subsetneq X$ be closed and put $U=X\setminus Y$.
Then $Q(X,Y)$ is topologically finitely generated if either
of the following conditions holds:
\begin{enumerate}
\item $\pi_1(U_{\overline\eta})$ is topologically finitely
generated;
\item $X_{\overline\eta}$ is normal and
$Q(X_{\overline\eta},Y_{\overline\eta})$ is topologically
finitely generated.
\end{enumerate}
The dimensions of $X$, $B$, and the fibres are unrestricted.
\end{proposition}
\begin{proof}
Write $F=k(B)$, $K=k(X)$, $G=Q(X,Y)$, and let $N$ be
the image of $\pi_1(U_{\overline\eta})$ in $G$.
The arithmetic fundamental-group exact sequence and
normality of $U$ show that $N$ is closed and normal,
and that $Q=G/N$ is a quotient of
$\operatorname{Gal}(F^{\mathrm{sep}}/F)$.
The proof is exactly the first paragraph of
Proposition~\ref{D:generic-fibre-generation}, using
\cite[Lemma 58.14.3, tag 0BTX]{C_StacksGenericFibre}.
In case (1), $N$ is finitely generated by hypothesis.
In case (2), restrictions of original tame covers to
the geometric generic fibre remain tame at its
codimension-one boundary DVRs, by
\cite[Lemma 15.116.9, tag 0EXY]{C_StacksLogFibration}.
Thus $N$ is a quotient of the finitely generated group
in (2), and is again finitely generated.

We obtain one base open that works for \emph{every}
finite quotient of $Q$.  The smooth locus of $U\to B$
dominates $B$, because the geometrically integral generic
fibre contains a dense smooth open
\cite[Lemma 33.25.7, tag 056V]{C_StacksGenericFibre}.
Let $B^\circ$ be its open image and put
$T=U_{\mathrm{sm}}\times_B B^\circ$.
Then $T\to B^\circ$ is smooth, surjective and
quasi-compact, hence fpqc.

For a finite quotient $Q_0$ of $Q$, its connected cover
$V\to U$ is generically the pullback of a finite Galois
extension $F'/F$, with group $Q_0$.  Let $B'$ be the
normalization of $B$ in $F'$.  Normalization commutes
with smooth base change
\cite[Lemma 37.19.2, tag 03GV]{C_StacksProjectiveBase}.
Consequently $B'\times_B T$ is the normalization of
$T$ in its generic cover, and is therefore $V\times_U T$.
The latter is finite \'etale over $T$.  \'Etaleness is
fpqc local on the target
\cite[Lemma 35.23.31, tag 02VN]{C_StacksProjectiveBase},
so $B'\to B$ is finite \'etale over $B^\circ$.
This proves, including at codimension-two points, that
the surjection from the absolute Galois group factors as
\[
                         \pi_1(B^\circ)\longrightarrow Q.
\]

There are finitely many prime divisors $D$ in
$B\setminus B^\circ$.  For each one choose a prime divisor
$E$ of $X$ dominating $D$.  Such a divisor exists by
properness and dominance: over the DVR at the generic
point of $D$, a nonzero uniformizer cuts a nonempty
principal divisor on the integral source.  Write $m_E$
for its positive valuation at $E$, and set
\[
                        F_D=\kappa(D),\qquad M_E=\kappa(E).
\]
The extension $M_E/F_D$ is finitely generated.  In positive
characteristic fix a finite purely inseparable field
extension $F_{D,0}/F_D$ from
Lemma~\ref{D:residual-inseparability-bound}, and put
$c_D=[F_{D,0}:F_D]$.  These choices are made once, before
any finite quotient is chosen.

Return to a finite quotient $Q_0$, with fields $F'$ and
$k(V)=KF'$.  Choose a prolongation of the $E$-valuation
to $KF'$, and restrict it to a valuation $D'$ of $F'$.
Let $e_E$ be its ramification index over $E$ and $e_D$
the index at $D'/D$.  Multiplication of value-group
indices gives
\[
                         e_D\mid m_Ee_E.
\tag{D.PB.1}
\]
If $E$ is retained in $U$, the cover is \'etale there;
otherwise $E$ is an original boundary divisor, and the
cover is tame there.  In either case $e_E$ is prime to
$p$ and the residual extension above $E$ is separable.
The latter contains the compositum of $M_E$ and
$L=\kappa(D')$.  Hence $LM_E/M_E$ is separable, and
Lemma~\ref{D:residual-inseparability-bound} gives
\[
       [\kappa(D'):F_D]_{\mathrm{insep}}\mid c_D,
       \qquad (e_D)_p\mid (m_E)_p.
\tag{D.PB.2}
\]

For clarity, the second bound alone would not control
fierce inertia.  For a finite Galois extension of these
divisorial DVRs, the full inertia order is
\[
        |I_{D'}|=e_D[\kappa(D'):F_D]_{\mathrm{insep}},
\]
and its wild subgroup has order
\[
        |P_{D'}|=(e_D)_p
                     [\kappa(D'):F_D]_{\mathrm{insep}}.
\]
These formulas include inseparable residue extensions.
The integral closures are finite; the degree formula and
the decomposition action on the normal residue extension
therefore give $|I|=e f_{\mathrm{insep}}$ without a defect
term.  The wild subgroup is a $p$-group and the tame
quotient has order the prime-to-$p$ part of $e$, by
\cite[Lemma 15.114.5, tag 09EE]{C_StacksProjectiveBase}.
Thus (D.PB.2) bounds the image of the absolute wild
inertia at $D$ in every finite quotient of $Q$ by the
fixed integer $(m_E)_p c_D$.  Its image $\overline P_D$
in $Q$ is finite.  This conclusion uses the bounded
finite quotient orders and the surjectivity of wild
inertia onto finite-level wild inertia; the latter is
included in \cite[Lemma 58.13.5, tag 0BUA]
{C_StacksProjectiveBase}.  In characteristic zero all
wild inertia groups are trivial.

Choose an open normal subgroup $Q'\subset Q$ meeting all
the finitely many finite groups $\overline P_D$ trivially.
Its inverse image in $\pi_1(B^\circ)$ defines a connected
finite \'etale cover $B_1^\circ\to B^\circ$.
The normalization $B_1$ of $B$ in its function field is
normal and projective, since it is finite over $B$.
Every codimension-one component of
$B_1\setminus B_1^\circ$ lies over a codimension-one
component of $B\setminus B^\circ$.  Its wild inertia
maps into a conjugate of the corresponding original
wild inertia: it is pro-$p$, maps into original inertia,
and that inertia modulo wild inertia is pro-prime-to-$p$.
Its image in $Q'$ is therefore trivial.  Consequently
\[
       \pi_1(B_1^\circ)\longrightarrow Q'
\]
factors through $Q(B_1,B_1\setminus B_1^\circ)$.
This last group is finitely generated by
Proposition~\ref{D:fixed-boundary-positive}, because
$B_1$ is projective.  Hence $Q'$, its finite extension
$Q$, and finally $G$ are finitely generated.
\end{proof}

\begin{corollary}[Proved: vertical boundary and low-dimensional generic fibres over projective images]
\label{D:projective-image-corollaries}
Let $X$ be a normal connected proper scheme over an
algebraically closed field, and suppose $X$ admits a
nonconstant morphism to a projective scheme.
Replace the target by its image and then by the normal
projective Stein factor $B$.
\begin{enumerate}
\item If $Y\subsetneq X$ is closed and its image is a
proper subset of $B$, then $Q(X,Y)$ is topologically
finitely generated.
\item If the normalization of the reduced geometric
generic fibre is projective, then $Q(X,Y)$ is
topologically finitely generated for every proper
closed subset $Y\subset X$.
\end{enumerate}
In particular (2) applies whenever the generic fibre
has dimension at most one.  In either assertion, if
$Y$ has codimension at least two, the group in question
is the full $\pi_1(X\setminus Y)$.
\end{corollary}
\begin{proof}
The function field $F=k(B)$ is algebraically closed in
$K=k(X)$, by Stein factorization and normality of $X$.
Consequently the geometric generic fibre is irreducible;
it need not be reduced when $\dim B>1$.
By \cite[Lemma 33.27.4, tag 0BXT]{C_StacksLogFibration},
choose a finite purely inseparable extension $F_0/F$
such that the normalization of the generic fibre after
base change is geometrically normal.  Normalize $B$
in $F_0$ and $X$ in $KF_0$ to obtain $X_0\to B_0$.
Then $B_0$ remains normal and projective, and $X_0\to X$
is finite and radicial, by the fixed-power integral
argument in Theorem~\ref{D:surface-map-all-boundaries}.
In characteristic zero we may take the trivial extension.

The geometric generic fibre of $X_0$ is the normalization
of the reduced original geometric generic fibre: the
finite birational map has normal source and hence is
its normalization.  It is geometrically integral.
In case (1) the geometric generic boundary is empty,
and this fibre is proper and connected.  Its full
fundamental group is finitely generated by
\cite[X, Theorem 2.9, pp. 273--274]{SGA1}.
In case (2) its fixed-boundary group is finitely generated
by the projective case of
Proposition~\ref{D:fixed-boundary-positive}.
The respective alternatives of
Proposition~\ref{D:projective-base-generation} therefore
give finite generation of $Q(X_0,Y_0)$.

Topological invariance of the \'etale site and tame
base change at the corresponding boundary DVRs give
a surjection $Q(X_0,Y_0)\to Q(X,Y)$, as proved in
Theorem~\ref{D:surface-map-all-boundaries}.  This proves
both assertions.  A normal proper curve over an
algebraically closed field is projective, as is a
zero-dimensional proper scheme, so the indicated
generic-fibre dimension bound implies (2).
Finally, in codimension at least two there are no
original boundary inertia groups to quotient out.
\end{proof}
